
\section{Signed inward transport in one unforced flow and the radial concentration barrier}
\label{sec:xl-inward-concentration}
The exact moving-cylinder conservation law of Section~\ref{sec:xxxix-moving}
does not establish the \emph{sign} of nonlinear transport. We first
construct a single smooth, globally regular, exactly unforced
three-dimensional flow with inward angular-momentum transport across
every sufficiently small radius for a common positive interval.
We then quantify the concentration and $H^2$ cost of a \emph{hypothetical}
divergent radial-flux amplitude. These results are not a construction
of unforced finite-time blowup.

Put $r^2=x^2+y^2$, $\rho^2=r^2+z^2$, $\varepsilon=\nu/1000$,
and fix the explicit nonnegative normalized axial cutoff
\[
\eta(z)=
\frac{e^{-1/(1-z^2)}\,\mathbf1_{\{|z|<1\}}}
{\int_{-1}^{1}e^{-1/(1-s^2)}ds},\qquad\int\eta(z)dz=1.
\]
The numerator is extended smoothly by zero.
For $r=R>0$ define the complete radial \emph{advective}
angular-momentum flux
\[
\mathcal F(R,t)=R^2\int_{\mathbb R}\int_0^{2\pi}
\eta(z)u_r(R,\theta,z,t)u_\theta(R,\theta,z,t)
\,d\theta\,dz.
\]
Thus $\mathcal F=F_R^{\rm adv}$ in Edition~XXXIX; neither
pressure nor viscosity is concealed in this term.

\begin{theorem}[A common-time, infinite-scale signed inward flux]
\label{thm:xl-single-flow}
Let the whole-space initial datum be
\begin{equation}\label{eq:xl-datum}
u_0=\varepsilon\operatorname{curl}
[e^{-\rho^2}(x-r^2y)e_z].
\end{equation}
There exists a unique global smooth finite-energy solution of
the \emph{exact unforced} three-dimensional Navier--Stokes equations
with this nonzero solenoidal Schwartz datum. With
$\Gamma=ru_\theta$, its initial angular-mean covariance,
torque and cylinder flux are
\begin{align}
\overline{u'_r\Gamma'}(0,r,z)
&=-\varepsilon^2r^3e^{-2(r^2+z^2)},\\
\mathcal T(0,r,z)
&=4\varepsilon^2r^2(1-r^2)e^{-2(r^2+z^2)},\\
\mathcal F(R,0)
&=-2\pi\varepsilon^2R^4e^{-2R^2}
\int\eta(z)e^{-2z^2}dz<0.
\end{align}
There exist $R_*>0$ and $\delta>0$ (not numerically
evaluated) such that the \emph{one fixed unforced flow} obeys
\begin{equation}\label{eq:xl-common-interval}
\boxed{\quad
\mathcal F(R,t)\le-\pi\varepsilon^2e^{-2}R^4
\qquad(0<R\le R_*,\ 0\le t\le\delta).\quad}
\end{equation}
In particular it has signed inward flux through every radius
$R_j=2^{-j}R_*$ simultaneously throughout this same time interval.
\end{theorem}
\begin{proof}
The datum is a smooth solenoidal Schwartz curl and is
nonaxisymmetric and reflection-breaking.
It differs from the Gaussian of Edition~XXXVII solely by
reversing the sign of the $r^2y$ potential mode.
Its unitary Fourier absolute-value majorant is unchanged,
so $C_F\|u_0\|_{\mathcal X^{-1}}/\nu
\le1/(100\sqrt\pi)<1/16$.
The explicitly contracting Fourier mild equation in
Edition~XXXIV therefore yields a global smooth unforced
solution with reconstructed physical pressure.

Let $f=e^{-(r^2+z^2)}$ and
$A_z=\varepsilon f(r\cos\theta-r^3\sin\theta)$.
Direct cylindrical differentiation gives
\[
u_{0,r}=\varepsilon f(-\sin\theta-r^2\cos\theta),\qquad
u_{0,\theta}=\varepsilon f[
(2r^2-1)\cos\theta-(2r^4-3r^2)\sin\theta].
\]
Angular integration gives the displayed covariance,
and the complete angularly averaged nonlinear advection
agrees with the stated torque
$-r^{-1}\partial_r(r\overline{u'_r\Gamma'})$.
The cylinder formula follows.

For every smooth Cartesian velocity near the $z$-axis,
the angular Taylor expansion shows that
$\overline{u'_r\Gamma'}=r^3a(r,z,t)$, with a continuous
$a$ through $r=0$, continuously depending on $z,t$
on compact sets. For the exact solution at $t=0$,
$a(0,z,0)=-\varepsilon^2e^{-2z^2}\le
-\varepsilon^2e^{-2}$ for $|z|\le1$.
Smooth time-continuity supplies common
$R_*,\delta>0$ with
$a(r,z,t)\le-\varepsilon^2e^{-2}/2$ throughout
$0\le r\le R_*$, $|z|\le1$, $0\le t\le\delta$.
The normalized nonnegative cutoff then yields
\eqref{eq:xl-common-interval}.
\end{proof}

\begin{proposition}[Positive angular-momentum accumulation under the full PDE]
\label{prop:xl-moment-production}
For the exact solution of Theorem~\ref{thm:xl-single-flow},
the complete weighted cylinder moment of Edition~XXXIX
satisfies, at every fixed $R>0$,
\begin{equation}\label{eq:xl-moment-production}
\boxed{\quad
\mathcal M'_{R,\eta}(0)
=-\mathcal F(R,0)
=2\pi\varepsilon^2R^4e^{-2R^2}
\int\eta(z)e^{-2z^2}dz>0.\quad}
\end{equation}
The identity remains true for an initially moving radius,
because its geometric sweep contribution is zero at $t=0$.
\end{proposition}
\begin{proof}
The mean initial angular momentum $\bar\Gamma(0,r,z)$
and all its derivatives vanish, and $u_z(0)=0$.
Thus the mean radial viscous, axial and moving-boundary
terms in the complete cylinder identity vanish at $t=0$.
The only remaining term is the explicitly signed
advective flux $-\mathcal F(R,0)$.
\end{proof}

\begin{theorem}[Energy localization of amplified radial-flux spikes]
\label{thm:xl-concentration}
Let $u$ be any smooth finite-energy exactly unforced
Navier--Stokes flow and $E_0=\|u(0)\|_2^2>0$.
Set $B_\eta=\frac12\|\eta\|_\infty E_0$.
At each time before any hypothetical singularity,
\begin{equation}\label{eq:xl-logbudget}
\boxed{\quad
\int_0^\infty\frac{|\mathcal F(R,t)|}{R}\,dR
\le B_\eta.\quad}
\end{equation}
Hence, if $|\mathcal F(R,t)|\ge A>0$ on a
radial interval $I$, its logarithmic width satisfies
$\int_I dR/R\le B_\eta/A$.
Moreover, if $|\mathcal F(R_0,t)|\ge A$,
some $R_c\in[R_0,R_0e^{4B_\eta/A}]$ obeys
\begin{equation}\label{eq:xl-steep}
\boxed{\quad
|R_c\partial_R\mathcal F(R_c,t)|
\ge \frac{3A^2}{16B_\eta}.\quad}
\end{equation}
In particular, a hypothetical terminal flux
$|\mathcal F(R_\tau,1-\tau)|\ge c\tau^{-h}$
with $h,c>0$ cannot remain comparably large
across a whole dyadic annulus.
Every interval on which its magnitude is at least
$c\tau^{-h}/2$ has logarithmic width at most
$2B_\eta c^{-1}\tau^h$.
\end{theorem}
\begin{proof}
Since $\Gamma=Ru_\theta$,
$2|u_ru_\theta|\le |u|^2$ and the unforced
energy inequality imply
\[
\int_0^\infty|\mathcal F(R,t)|\,dR/R
\le\tfrac12\int_{\mathbb R^3}|\eta(z)||u(x,t)|^2dx
\le B_\eta.
\]
This proves the interval bound.
On a logarithmic interval of length $4B_\eta/A$
beginning at $\log R_0$, the average of
$|\mathcal F|$ is at most $A/4$; choose a radius
with that value and apply the mean-value theorem
to $\mathcal F(e^s,t)$.
The half-height interval estimate follows directly.
\end{proof}

\begin{proposition}[Necessary high-derivative cost of a divergent flux]
\label{prop:xl-h2-flux}
With the explicit whole-space interpolation constant
$C_2=4/(\pi\sqrt2\,3^{7/8})$ from Edition~XXXIII,
\begin{equation}\label{eq:xl-h2}
\boxed{\quad
\|\Delta u(t)\|_2\ge
\left(
\frac{|\mathcal F(R,t)|}
{\pi R^2\|\eta\|_1C_2^2E_0^{1/4}}
\right)^{2/3}.\quad}
\end{equation}
If $R_\tau=\sqrt{2X\tau}$ and a hypothesis
$|\mathcal F(R_\tau,1-\tau)|\ge c\tau^{-h}$
holds for all sufficiently small $\tau>0$,
then
\[
\|\Delta u(1-\tau)\|_2\ge
\left(
\frac{c}{2\pi X\|\eta\|_1C_2^2E_0^{1/4}}
\right)^{2/3}\tau^{-2(1+h)/3},
\]
so $\int_{t_0}^1\|\Delta u(t)\|_2^2dt=\infty$.
\end{proposition}
\begin{proof}
The flux definition gives
$|\mathcal F(R,t)|\le
\pi R^2\|\eta\|_1\|u(t)\|_\infty^2$.
Combine this with
$\|u\|_\infty\le C_2\|u\|_2^{1/4}
\|\Delta u\|_2^{3/4}$ and $\|u(t)\|_2^2\le E_0$.
The squared terminal lower exponent
$4(1+h)/3>1$ is nonintegrable.
\end{proof}

\paragraph{Scope.}
The first theorem proves inward signed flux at every small
radius over a common initial interval of \emph{one global
regular unforced flow}. It does not prove time-dependent
terminal amplification. The later estimates constrain
the distinct, stronger hypothesis that the \emph{radial
flux amplitude itself} grows as $\tau^{-h}$.
The preceding source-specific escape requirement concerns
the time-integrated torque or angular variation and
does not imply that radial-flux hypothesis.
No actual OpenAI Section~10 numerical source
profile, singularity-preserving zero-force correction,
or complete Lean PDE certificate has been obtained.
